% =====================================================================
%  27.3  Mo on PATHS -- the doubled Euler tour E[1..2n] with the
%  membership TOGGLE, and the lca split (ancestor case vs general case,
%  the LCA-parity trap).  WORKED EXAMPLES ONLY -- no solvedproblem here
%  (chapter contract: exactly two, in 27.1 and 27.2).
%  Exports: sec:c21-path  def:c21-toggle  thm:c21-pathsplit
%  lst:c21-euler-toggle
%  Every trace below was replayed by hand against the authoritative
%  running example of 00_opener: E=[1,2,4,8,8,4,5,5,2,3,6,6,7,7,3,1],
%  st(v) node-ordered = (1,2,10,3,7,11,13,4), en(v) = (16,9,15,6,8,12,14,5),
%  val = (2,1,3,1,2,3,1,2), lca from Ch.16 usage (never re-derived).
%  CF 258E / 86D style ideas are PROSE-ONLY; their URLs live elsewhere.
% =====================================================================
\section{Mo on Paths: the Doubled Tour and the Toggle}\label{sec:c21-path}

\begin{previouslybox}[27.2 made subtrees ranges -- paths refuse that]
The flatten of \secref{c21-subtree} works because subtree\((v)\) is a
\emph{contiguous} block of preorder.  A path \(u\)-\(v\) is not: its nodes
\(u, \parent(u), \dots, \lca, \dots, \parent(v), v\) climb out of one subtree
and down another, and no re-ordering of a length-\(n\) array makes every one of
\(\bigO{n^2}\) candidate paths contiguous at once.  The escape is to stop
demanding "the range \emph{equals} the path" and settle for "the range
\emph{represents} the path" -- via the doubled Euler tour \(E\) of the opener,
where each node is written at entry and at exit, and membership is read off by
\textbf{parity of occurrence count}.
\end{previouslybox}

\begin{topicintro}
Length \(2n\) tour, two occurrences per node.  As a Mo pointer crosses a
position it \term{toggles} that node: appears once in the window means
"\emph{in} the represented set", twice means "\emph{out}".  \thmref{c21-pathsplit}
says exactly which window represents path \(u\)-\(v\): order so
\(st(u)<st(v)\); if \(u\) is an ancestor of \(v\) the window \([st(u),st(v)]\)
represents the whole path; otherwise \([en(u),st(v)]\) represents the path
\emph{minus} \(\lca(u,v)\) -- and the LCA must be fed to the statistic by hand,
counted once, then withdrawn.  That last clause is the famous parity trap, and
this section disarms it with three hand-replayed walks on the running tree \(R\):
path\((4,8)\) (ancestor case), path\((8,6)\) (general case), path\((2,3)\) (the
trap \emph{changing an answer}).

\textbf{Prerequisites.} \secref{c21-moengine} (window + \texttt{add}/\texttt{rm}),
\secref{c21-subtree} (the flatten mindset), \(\lca\) as a black box from
\chref{c11} (\(\thmref{c11-distf}\)'s chapter -- Tarjan gets its own offline
implementation in \secref{c21-tarjan}).

\textbf{After this section you can}
\begin{itemize}[nosep]
  \item build the doubled tour and toggle-window it in your sleep; \hfill\lv{4}
  \item state both windows of \thmref{c21-pathsplit} and prove the parity claim; \hfill\lv{5}
  \item handle the ancestor case and the add-the-LCA-by-hand case correctly; \hfill\lv{5}
  \item answer distinct-colour-on-path offline in \(\bigO{(n+q)\sqrt n}\). \hfill\lv{5}
\end{itemize}
\end{topicintro}

% ---------------------------------------------------------------
\subsection{The doubled tour and the toggle window}
% ---------------------------------------------------------------

\begin{definition}{The doubled Euler tour and the toggle}{c21-toggle}
Run the DFS writing \(v\) once when you \emph{enter} it and once when you
\emph{leave} it: \(E[1..2n]\), with every node appearing exactly twice, at
\(st(v)\) (first) and \(en(v)\) (second).  On the running tree \(R\):
\[
  E=[1,2,4,8,8,4,5,5,2,3,6,6,7,7,3,1],
\]
\(st(v)=(1,2,10,3,7,11,13,4)\) and \(en(v)=(16,9,15,6,8,12,14,5)\) for
\(v=1..8\).  A Mo window \([\ell,r]\) over \(E\) \textbf{toggles} node membership
as pointers cross occurrences: a node is \emph{represented} iff it occurs an
odd number of times (i.e.\ exactly once) in \(E[\ell..r]\).  \texttt{add}/\texttt{rm}
are the same test: \(\mathit{held}[x] \oplus= 1\); insert \(x\) into the
statistic when \(\mathit{held}\) turns \(1\), erase when it turns \(0\).
\end{definition}

\begin{theorem}{Two windows represent a path (the LCA split)}{c21-pathsplit}
Let \(st(u)<st(v)\).  (i) If \(\lca(u,v)=u\), the odd-occurrence set of
\(E[st(u)\,..\,st(v)]\) is exactly the node set of path \(u\)-\(v\).
(ii) Otherwise the odd-occurrence set of \(E[en(u)\,..\,st(v)]\) is the path
node set \emph{minus} \(\lca(u,v)\).
\end{theorem}
\begin{proofbox}
A node occurs oddly in a window iff exactly one of \(st,\ en\) lands inside it.
\emph{Case (i):} window \([st(u),st(v)]\), \(u\) an ancestor of \(v\).  Any
path node \(w\) (including \(u\) and \(v\)): \(st(w)\in[st(u),st(v)]\) since
\(w\) lies on the descent from \(u\) to \(v\), while \(en(w)\ge en(v)>st(v)\)
since \(w\) ancestors \(v\) -- exactly one occurrence inside.  A node \(y\)
outside subtree\((u)\): both occurrences left of \(st(u)\) or both right of
\(st(v)\).  A node \(y\in\subtree(u)\) off the path: \(y\) hangs in a side
branch of some path node \(w\), so its span \([st(y),en(y)]\) is disjoint from
\(st(v)\) -- either entirely inside \([st(u),st(v)]\) (two occurrences, even)
or entirely to its right (zero).  Odd set \(=\) path.  \emph{Case (ii):}
\(w=\lca(u,v)\neq u\), window \([en(u),\ st(v)]\).  Ancestor \(x\) of \(u\)
strictly below \(w\): \(st(x)<st(u)<en(u)\) (outside, left) and
\(en(u)<en(x)<st(v)\) (inside: \(x\)'s subtree closes after \(u\)'s does, and
since \(v\notin\subtree(x)\) while \(st(x)<st(v)\), branch \(x\) finishes
before \(v\) opens) -- once.  Symmetrically, ancestor \(y\) of \(v\) strictly
below \(w\): \(st(y)\in(en(u),st(v))\) with \(en(y)>st(v)\) -- once.  \(u\):
only \(en(u)\) in the window; \(v\): only \(st(v)\).  \(w\) itself:
\(st(w)<st(u)\) and \(en(w)>st(v)\) -- \emph{zero} occurrences, the missing
endpoint.  Every other node has both occurrences inside one subtree span
entirely left of \(en(u)\), entirely right of \(st(v)\), or nested and fully
contained in the window -- always an even count.  Odd set
\(=\) path \(\setminus\{w\}\). \(\Box\)
\end{proofbox}

\begin{lstlisting}[caption={Doubled-tour build, toggle add/rm, and the window setup for a path query. The lca node is fed to the statistic between answering steps in case (ii). ASCII only, no underscores}, label={lst:c21-euler-toggle}]
int eu[2*MAXN], ecnt, st[MAXN], en[MAXN], held[MAXN];
void dfs(int u, int p){
    st[u] = ++ecnt;  eu[ecnt] = u;
    for (int v : adj[u]) if (v != p) dfs(v, u);
    en[u] = ++ecnt;  eu[ecnt] = u;          // second write = exit
}
void add(int pos){                            // toggle node at tour position
    int x = eu[pos];
    held[x] ^= 1;
    if (held[x]) insert(x);                   // entered the represented set
    else       erase(x);                      // left it
}
void rm(int pos){ add(pos); }                 // toggle is its own inverse
// window for query (u,v):
int L, R; bool extraLca;
if (st[u] > st[v]) swap(u, v);                // keep st[u] < st[v]
if (is_ancestor(u, v)){ L = st[u]; R = st[v]; extraLca = false; }
else                  { L = en[u]; R = st[v]; extraLca = true;  }
// is_ancestor(a,b): tin/tout from the SAME dfs, tin[a]<=tin[b] && tout[b]<=tout[a]
// (or st[a]<=st[b] && en[b]<=en[a] on this tour).
// In case extraLca:  answer = statistic(window UNION {lca(u,v)})
//                    -- call insert(lca), read cur, erase(lca).
\end{lstlisting}

\begin{insightbox}[Why toggle is legal where insert is not]
On subtrees, each node sits at \emph{one} array position, so "in the window"
was unambiguous.  On the doubled tour a node can cross the pointer twice --
a plain \texttt{add} would count it twice and a distinct-count would lie.
Making \texttt{add} \(\equiv\) \texttt{rm} (an xor-toggle) means two crossings
annihilate: the window's multiset of \emph{distinct represented nodes} always
equals the odd-occurrence set, no matter how the pointers wandered.  This is
exactly why path Mo tolerates the same \(O(1)\) block-order sweep as array Mo
(\(\thmref{c21-complexity}\) applies verbatim with \(n \to 2n\)).
\end{insightbox}

% ---------------------------------------------------------------
\subsection{Worked example: distinct colours on a path of \(R\)}
% ---------------------------------------------------------------

The statistic: \(\mathit{distinct}=\) number of different \(\mathit{val}\)
colours among represented nodes, maintained the classic way -- a colour
frequency array \(\mathit{cf}\), a counter of colours with \(\mathit{cf}>0\);
\texttt{insert(x)}: \(\mathit{cf}[\mathit{val}(x)]{+}{+}\) (bumping \(0\to1\)
raises the count); \texttt{erase(x)} mirrors it.  Colours on \(R\):
\(\mathit{val}=(2,1,3,1,2,3,1,2)\) for nodes \(1..8\).

\begin{tracebox}[Walk 1 -- ancestor case, path\((4,8)\)]
\(st(4)=3 < st(8)=4\); \(\lca(4,8)=4=u\) \(\Rightarrow\) case (i), window
\([3,4]\).  \(E[3..4]=[4,8]\): both occur once \(\Rightarrow\) represented
\(\{4,8\}\) -- the path \(4\to8\) exactly, LCA \(=\) endpoint \(4\) already
inside.  Colours \(\mathit{val}(4)=1,\mathit{val}(8)=2\): \(\boxed{2}\)
distinct.  No extra handling. \(\Box\)
\end{tracebox}

\begin{tracebox}[Walk 2 -- general case, path\((8,6)\)]
\(st(8)=4 < st(6)=11\); \(\lca(8,6)=1\neq8\) \(\Rightarrow\) case (ii), window
\([en(8)=5,\ 11]\).  \(E[5..11]=[8,4,5,5,2,3,6]\).  Occurrence parity:
\(8{:}1\), \(4{:}1\), \(5{:}2\) (toggles out), \(2{:}1\), \(3{:}1\), \(6{:}1\)
\(\Rightarrow\) represented \(\{8,4,2,3,6\}\).  The true path is
\(8\to4\to2\to1\to3\to6\) -- \(\lca=1\) is missing, precisely as
\thmref{c21-pathsplit}(ii) promises.  Feed node \(1\) to \texttt{insert},
read, then \texttt{erase}: represented \(\{8,4,2,1,3,6\}\), colours
\(2,1,1,2,3,3\) \(\Rightarrow\) \(\boxed{3}\) distinct. \(\Box\)
\end{tracebox}

\begin{tracebox}[Walk 3 -- the trap bites: path\((2,3)\) with and without the LCA]
\(st(2)=2 < st(3)=10\); \(\lca(2,3)=1\neq2\) \(\Rightarrow\) window
\([en(2)=9,\ 10]\): \(E[9..10]=[2,3]\), both odd \(\Rightarrow\) represented
\(\{2,3\}\), colours \(1,3\): \(\mathit{distinct}=2\).  But the path is
\(2\to1\to3\) and node \(1\) brings colour \(2\) \(\Rightarrow\) true answer
\(\boxed{3}\).  Forgetting the extra \(\mathit{lca}\) \emph{changes the
answer} -- on walk~2 it happened not to (colour \(2\) was already carried by
node~8), on walk~3 it does.  A single missed trace of this shape is a WA on
CF path-Mo problems. \(\Box\)
\end{tracebox}

\begin{complexitybox}[Path Mo]
Tour build \(\bigO{n}\); each query's window computed in \(O(1)\) given
\(\lca\) (Ch.~16 lifting: \(\bigO{q\log n}\) total) or precomputed batch.
Sweep: \(\bigO{(2n+q)\sqrt{2n}} = \bigO{(n+q)\sqrt n}\) toggles, each with the
\(O(1)\) inner \texttt{insert}/\texttt{erase}.  The add/read/erase of the LCA
is \(O(1)\) per query. \textbf{Space} \(\bigO{n}\) tour + statistic
arrays.  The same skeleton answers "how many nodes with \(\mathit{val}\ge X\)",
"\(\sum \mathit{cf}^2\)"-style pair counts, mode-of-colours on a path --
anything the toggle can maintain at \(O(1)\).
\end{complexitybox}

\begin{warningbox}[The three path-Mo failure modes]
\begin{enumerate}[nosep,leftmargin=1.6em]
  \item \textbf{Forgetting case (ii)'s LCA} (walk~3).  Symptom: answers wrong
        exactly when the LCA's colour is new to the window.  Fix: wrap the
        read with \texttt{insert(lca)}/\texttt{erase(lca)} -- and \emph{only at
        read time}, never as a permanent window member.
  \item \textbf{Swapping endpoints wrongly.} The rules assume \(st(u)<st(v)\);
        compare \(st\), not node id.  Then case (i) is "\(\lca = u\)", the
        \emph{left} endpoint -- if your code tests the wrong endpoint the
        ancestor case silently degrades to a wrong range.
  \item \textbf{Double-adding within case (i).} When \(u\) is the ancestor,
        \(u\) is \emph{already} oddly represented; inserting the LCA again
        toggles it \emph{out}.  The extra insert belongs to case (ii) only.
\end{enumerate}
\end{warningbox}

\begin{hardeningbox}
\begin{itemize}[nosep,leftmargin=1.4em]
  \item \textbf{Query sorting uses the window endpoints} \((L_i/B, R_i)\)
        computed \emph{after} the split -- sorting by raw \((st,r)\) mixes the
        two cases and breaks nothing, but the standard form keeps every window
        a plain range over \(E[1..2n]\).
  \item \textbf{Block size on \(2n\):} \(B \approx 2n/\sqrt q\); with
        \(q=\Theta(n)\), \(B\approx\sqrt{2n}\).
  \item \textbf{The toggle pattern is the tree analogue of "xor is its own
        inverse"}: \texttt{rm = add} because the second occurrence undoes the
        first.  Problems like CF~86D's \(\sum\mathrm{freq}^2\) on arrays do not
        need it; paths cannot live without it.  (86D's URL lives in
        \(\secref{c21-boundary}\)'s prose -- no re-linking.)
\end{itemize}
\end{hardeningbox}

\begin{recallbox}[27.3 in one breath]
\emph{Double the tour, \term{toggle} on every crossing, and a path becomes a
range: \([st(u),st(v)]\) when \(u\) ancestors \(v\), else \([en(u),st(v)]\)
\emph{plus the LCA handed to the statistic at read time}
(\(\thmref{c21-pathsplit}\), \(\coderef{c21-euler-toggle}\)).  Walks on \(R\):
\((4,8)\to2\), \((8,6)\to3\) with the LCA, \((2,3)\to3\) vs.\ a lying \(2\)
without it.  Cost stays \(\bigO{(n+q)\sqrt n}\).}
\selftest{
  \item Why must the represented set be defined by \emph{odd} occurrence? \emph{(Two crossings of the same node must annihilate; parity is the only toggle-invariant that makes rm = add.)}
  \item Give the window for path\((5,6)\) of \(R\) and its odd set. \emph{(\([en(5)=8,\ st(6)=11]=\) E[8..11] \(=[5,2,3,6]\), odd set \(\{5,2,3,6\}\), plus lca \(1\) at read time -- the path \(5,2,1,3,6\).)}
  \item When is the extra LCA insert \emph{not} needed? \emph{(Case (i): lca equals the left endpoint and is already oddly represented.)}
}
\end{recallbox}
